\mychapter{Dane i model obliczeniowy}
\label{chap:dane_i_model}

\section{Źródło danych: OpenStreetMap}
\label{sec:dane:osm}

Zarówno sieć transportowa, jak i lokalizacje usług (punkty zainteresowania, ang.\ \emph{Points of Interest}, PoI)
wykorzystywane w~niniejszej pracy pochodzą z~projektu OpenStreetMap (OSM) \cite{openstreetmap}.
OSM jest otwartą, ogólnoświatową bazą danych geoprzestrzennych, tworzoną i~utrzymywaną przez społeczność wolontariuszy.
Wybór tego źródła wynika z~trzech czynników:
\begin{itemize}
  \item \textbf{otwartość danych}: dane OSM są dostępne na licencji ODbL,
        co umożliwia ich swobodne wykorzystanie w~badaniach naukowych bez ograniczeń licencyjnych,
  \item \textbf{globalna replikowalność}: jednolity format danych i~schemat tagów pozwalają
        stosować te same narzędzia i~metody niezależnie od analizowanego miasta,
        co jest istotne dla powtarzalności eksperymentów,
  \item \textbf{jednolity schemat tagów}: OSM korzysta z~ustandaryzowanego systemu tagów
        (par klucz--wartość), który umożliwia systematyczne mapowanie kategorii usług
        na odpowiednie obiekty w~bazie danych.
\end{itemize}
Podobne uzasadnienie przyjęto m.in.\ w~pracy Olivariego i~in.\ \cite{olivari2023presenting}
(indeks NEXI) oraz w~pracy Lama \cite{Lam2024}, w~których OSM stanowił jedyne źródło danych wejściowych.

Dostęp do danych OSM realizowany jest za pośrednictwem biblioteki OSMnx \cite{boeing2017osmnx},
udostępniającej interfejs programistyczny w~języku Python do usługi Overpass API.
OSMnx automatyzuje pobieranie, filtrowanie i~konstrukcję grafów transportowych w~postaci obiektów
biblioteki NetworkX, dzięki czemu można bezpośrednio wykorzystać standardowe algorytmy grafowe.

\clearpage

\section{Konstrukcja grafu pieszego}
\label{sec:dane:graf}

Na podstawie danych OSM budowany jest graf ważony $G=(V,E)$ reprezentujący sieć pieszą badanego obszaru.
Proces konstrukcji obejmuje pięć kroków, podsumowanych poniżej.

\subsection{Pobranie danych}

Dane sieci pieszej pobierane są za pomocą jednej z~dwóch funkcji biblioteki OSMnx,
zależnie od sposobu zdefiniowania obszaru.
Dla obszarów identyfikowanych przez nazwę administracyjną (np.\ Gdańsk, Londyn)
wykorzystywana jest funkcja \texttt{graph\_\allowbreak from\_\allowbreak place},
która sama pobiera granice z~OSM i~konstruuje graf w~ich obrębie.
Dla obszarów złożonych z~kilku jednostek pomocniczych (np.\ Gdańsk Południe)
stosowana jest funkcja \texttt{graph\_\allowbreak from\_\allowbreak polygon},
której przekazywany jest wcześniej przygotowany wielokąt obszaru.
Wielokąt ten jest tworzony przez pobranie granic poszczególnych jednostek z~OSM
(funkcja \texttt{geocode\_\allowbreak to\_\allowbreak gdf})
i~połączenie ich operacją \texttt{unary\_\allowbreak union} z~biblioteki Shapely.
W~obu przypadkach stosowany jest parametr \texttt{network\_\allowbreak type="walk"},
który ogranicza wynik do ulic i~ciągów dozwolonych dla ruchu pieszego.

\subsection{Odwzorowanie kartograficzne}

Pobrany graf jest rzutowany do układu współrzędnych EPSG:3035 (metrycznego
odwzorowania równopolowego obejmującego Europę) za pomocą funkcji
\texttt{osmnx.\allowbreak projection.\allowbreak project\_\allowbreak graph}. Dzięki temu wszystkie odległości i~długości
krawędzi są wyrażone w~metrach, co jest konieczne do poprawnego wyznaczania wag czasowych.

\subsection{Konsolidacja skrzyżowań}

Surowe dane OSM często zawierają nadmiarowe węzły w~obrębie skrzyżowań,
np.\ osobne węzły dla każdego pasa ruchu lub przejścia.
Konsolidacja skrzyżowań, realizowana funkcją \texttt{osmnx.\allowbreak consolidate\_\allowbreak intersections}
z~parametrem \texttt{tolerance=10\,m}, łączy w~pojedyncze wierzchołki te węzły, które
leżą bliżej siebie niż zadana tolerancja. Przyjęta wartość $10\,\text{m}$ stanowi kompromis
między redukcją nadmiarowych węzłów a~zachowaniem szczegółów topologicznych sieci.

\subsection{Konwersja do grafu nieskierowanego}

Sieć piesza jest symetryczna (piesi mogą poruszać się w~obu kierunkach każdym odcinkiem),
dlatego wynikowy graf konwertowany jest do postaci nieskierowanej za pomocą funkcji
\texttt{osmnx.convert.to\_undirected}.

\subsection{Geometria krawędzi}

Krawędzie grafu mogą posiadać atrybut \texttt{geometry} przechowujący obiekt \texttt{LineString},
tj.\ łamaną złożoną z~wielu punktów pośrednich, odwzorowującą rzeczywisty przebieg odcinka ulicy.
Długość takiej krawędzi jest sumą odległości między kolejnymi punktami łamanej.
Krawędzie pozbawione tego atrybutu traktowane są jako odcinki prostoliniowe łączące współrzędne
wierzchołków końcowych, a~ich długość przyjmuje się równą odległości euklidesowej między
tymi wierzchołkami.

\section{Kategorie usług i mapowanie tagów OSM}
\label{sec:dane:uslugi}

W~niniejszej pracy wykorzystywany jest zestaw ośmiu kategorii usług zdefiniowanych
w~indeksie NEXI (Olivari i~in.\ \cite{olivari2023presenting}, Table~1).
Każda kategoria jest mapowana na odpowiedni zestaw tagów OSM,
co umożliwia automatyczne pozyskiwanie lokalizacji usług z~bazy danych.
Mapowanie przedstawiono w~tabeli~\ref{tab:nexi_tags}.

\begin{table}[ht]
\centering
\caption{Kategorie usług NEXI i ich mapowanie na tagi OSM.}
\label{tab:nexi_tags}
\begin{tabular}{l l p{7.5cm}}
\hline
\textbf{Kategoria} & \textbf{Klucz OSM} & \textbf{Przykładowe wartości} \\
\hline
education      & \texttt{amenity}  & college, kindergarten, school, university, \ldots \\
entertainment  & \texttt{amenity}  & cinema, theatre, community\_centre, arts\_centre, \ldots \\
grocery        & \texttt{shop}     & bakery, supermarket, convenience, greengrocer, \ldots \\
health         & \texttt{amenity}  & clinic, doctors, hospital, pharmacy, \ldots \\
posts\_banks   & \texttt{amenity}  & atm, bank, post\_office, bureau\_de\_change \\
parks          & \texttt{leisure}  & park, dog\_park \\
sustenance     & \texttt{amenity}  & restaurant, cafe, bar, fast\_food, pub, \ldots \\
shops          & \texttt{shop}     & clothes, electronics, furniture, hairdresser, \ldots \\
\hline
\end{tabular}
\end{table}

Przy mapowaniu tagów istotne są trzy decyzje projektowe:
\begin{itemize}
  \item \textbf{kategorie \emph{grocery} i~\emph{shops}}: obie wykorzystują klucz \texttt{shop}.
        Kategoria \emph{grocery} obejmuje sklepy spożywcze i~napojowe
        (np.\ bakery, supermarket, convenience), natomiast \emph{shops} obejmuje handel
        nieżywnościowy (np.\ clothes, electronics, furniture).
        Cztery wartości ogólnohandlowe (\texttt{department\_store}, \texttt{general},
        \texttt{kiosk}, \texttt{mall}) występują w~obu listach zgodnie z~oryginałem~\cite{olivari2023presenting},

  \item \textbf{kategoria \emph{parks}}: w~tabeli NEXI jest ona konceptualnie grupowana
        pod kluczem \texttt{amenity}, jednak w~schemacie tagów OSM wartości \texttt{park}
        i~\texttt{dog\_park} należą do klucza \texttt{leisure}. Zapytanie \texttt{amenity=park}
        zwraca w~OSM niemal puste wyniki; poprawne jest \texttt{leisure=park},

  \item \textbf{normalizacja geometrii}: obiekty pobrane z~OSM mogą mieć zróżnicowaną geometrię,
        np.\ punkty, linie lub wielokąty (obrys budynku szpitala).
        Wszystkie geometrie sprowadzane są do centroidów obliczanych w~układzie EPSG:3035,
        co daje jednolity format punktowy wymagany przez procedurę wstawiania usług do grafu.
\end{itemize}

\section{Wstawianie lokalizacji usług do grafu}
\label{sec:dane:wstawianie}

Aby umożliwić obliczanie odległości grafowych między wierzchołkami sieci pieszej
a~lokalizacjami usług, punkty usług muszą zostać wstawione do grafu jako dodatkowe wierzchołki.
Procedura wstawiania składa się z~następujących kroków, wykonywanych dla każdego punktu usługi:

\begin{enumerate}
  \item \textbf{znalezienie najbliższej krawędzi}: dla punktu usługi o~współrzędnych
        $(x_s, y_s)$ wyznaczana jest najbliższa krawędź $(u,v)$ grafu za pomocą funkcji
        \texttt{osmnx.\allowbreak distance.\allowbreak nearest\_\allowbreak edges},

  \item \textbf{najbliższy punkt na krawędzi}: współrzędne punktu rzutu $(x_s, y_s)$
        na geometrię krawędzi wyznacza się dwoma kolejnymi wywołaniami Shapely dla tej samej
        łamanej: \texttt{LineString.project} oraz \texttt{LineString.interpolate}.
        Metoda \texttt{project} zwraca liczbę $d$ równą długości fragmentu łamanej
        od jej początku do punktu na niej leżącego najbliżej $(x_s, y_s)$.
        Metoda \texttt{interpolate($d$)} odwzorowuje taką odległość na punkt $p$ leżący
        na łamanej, którego współrzędne są używane dalej w~procedurze.

  \item \textbf{podział krawędzi}: oryginalna krawędź $(u,v)$ jest usuwana,
        a~w~jej miejsce wstawiana jest sekwencja nowych krawędzi przechodzących przez
        punkty rzutów. Jeśli na tę samą krawędź rzutuje $k$ punktów usług,
        zostaje ona dzielona za jednym razem na $k{+}1$ segmentów (w~przypadku jednej
        usługi otrzymujemy dwie krawędzie: $(u, p)$ i~$(p, v)$). Geometrie nowych krawędzi
        wyznacza funkcja \texttt{shapely.ops.substring}, zwracająca podciąg łamanej
        między zadanymi odległościami,

  \item \textbf{utworzenie wierzchołka usługi}: w~każdym punkcie rzutu tworzony jest
        nowy wierzchołek grafu, oznaczony kategorią usługi (np.\ \emph{grocery}, \emph{health}),

  \item \textbf{wyznaczenie wag}: krawędziom grafu przypisywane są wagi czasowe zgodnie
        z~modelem opisanym w~podrozdziale~\ref{sec:dane:wagi}. W~implementacji wagi
        przeliczane są dla wszystkich krawędzi grafu, zarówno nowo powstałych,
        jak i~niezmienionych przez procedurę wstawiania, co zapewnia jednolitość modelu
        czasowego.
\end{enumerate}

W~celu poprawy wydajności najbliższe krawędzie dla wszystkich punktów usług
wyznaczane są jednym wywołaniem funkcji \texttt{osmnx.\allowbreak distance.\allowbreak nearest\_\allowbreak edges},
zamiast wielokrotnego wywoływania jej w~pętli, raz na każdy punkt.
Następnie punkty grupowane są według wspólnej krawędzi tak, aby każda krawędź była dzielona
tylko raz, niezależnie od liczby rzutowanych na nią punktów.

Podejście to jest analogiczne do metody zastosowanej przez Lama
(Figure~6.2 w~\cite{Lam2024}).

\section{Model wag czasowych}
\label{sec:dane:wagi}

Waga krawędzi $e \in E$ jest interpretowana jako czas przejścia pieszo
i~wyznaczana według wzoru
\[
w(e) = \frac{\ell(e)}{v},
\qquad
v = 5\,\text{km/h}
\approx 83{,}33\,\text{m/min}
\]
gdzie $\ell(e)$ oznacza długość geometrii krawędzi $e$ w~metrach,
a~$v$ przyjętą prędkość marszu wyrażoną w~metrach na minutę.
Tak zdefiniowana waga ma jednostkę \emph{minut}, co jest spójne
z~kontekstem miasta 15-minutowego, gdzie analizowany próg dostępności
($15\,$min) wyrażony jest w~tych samych jednostkach.
Wartość prędkości $5\,\text{km/h}$ jest zgodna z~założeniem przyjętym
w~pracy definiującej indeks NEXI \cite{olivari2023presenting}.

Model wag jest świadomie uproszczony: nie uwzględnia dodatkowych składników opóźnień
(np.\ na skrzyżowaniach czy przejściach dla pieszych) ani zróżnicowania prędkości
ze względu na typ nawierzchni lub nachylenie terenu.

\section{Obszary badawcze i charakterystyka grafów}
\label{sec:dane:obszary}

Eksperymenty prowadzone są na trzech obszarach o~zróżnicowanej skali,
co pozwala ocenić zachowanie badanych algorytmów w~warunkach obejmujących
zarówno analizę lokalną, jak i~test na wielkiej skali.

\subsection{Gdańsk Południe}

Gdańsk Południe obejmuje cztery jednostki pomocnicze miasta Gdańska:
Chełm, Jasień, Orunia Górna--Gdańsk Południe oraz Ujeścisko-Łostowice \cite{wiki:gdansk_poludnie}.
Granice tych jednostek są pozyskiwane z~OpenStreetMap za pomocą geokodowania,
a~następnie łączone w~jeden wielokąt operacją \texttt{unary\_union}.
Jest to najmniejszy z~badanych obszarów, stanowiący przykład analizy na poziomie dzielnicy.

\subsection{Gdańsk}

Obszar obejmuje całe miasto Gdańsk w~granicach administracyjnych.
Granica jest pozyskiwana z~OSM jako pojedynczy wielokąt odpowiadający relacji administracyjnej miasta.

\subsection{Londyn (Greater London)}

Greater London stanowi największy z~badanych obszarów i~pełni rolę testu skalowalności.
Granica jest pozyskiwana z~OSM analogicznie jak w~przypadku Gdańska.

\subsection{Wizualizacje grafów}

Rysunki~\ref{fig:graf_gdansk_poludnie}--\ref{fig:graf_london} przedstawiają grafy piesze
trzech badanych obszarów po wstawieniu lokalizacji usług.
Szare punkty odpowiadają wierzchołkom sieci pieszej (skrzyżowaniom),
natomiast czerwone oznaczają wstawione lokalizacje usług.

\begin{figure}[ht]
\centering
\includegraphics[width=0.75\textwidth]{gdansk_poludnie_walk_services.png}
\caption{Graf pieszy Gdańska Południe z~lokalizacjami usług.}
\label{fig:graf_gdansk_poludnie}
\end{figure}

\begin{figure}[ht]
\centering
\includegraphics[width=0.75\textwidth]{gdansk_walk_services.png}
\caption{Graf pieszy Gdańska z~lokalizacjami usług.}
\label{fig:graf_gdansk}
\end{figure}

\begin{figure}[ht]
\centering
\includegraphics[width=0.75\textwidth]{london_walk_services.png}
\caption{Graf pieszy Londynu z~lokalizacjami usług.}
\label{fig:graf_london}
\end{figure}

\clearpage

\subsection{Charakterystyka grafów}

Tabela~\ref{tab:grafy} przedstawia rozmiary grafów pieszych
dla trzech badanych obszarów (po konsolidacji skrzyżowań, przed wstawieniem wierzchołków usług).
Rozmiary różnią się o~ponad rząd wielkości,
co pozwala ocenić skalowalność algorytmów w~szerokim zakresie.

\begin{table}[ht]
\centering
\caption{Rozmiary grafów pieszych dla badanych obszarów
         (po konsolidacji skrzyżowań, przed wstawieniem usług).}
\label{tab:grafy}
\begin{tabular}{l r r r}
\hline
\textbf{Obszar} & $|V|$ & $|E|$ & \textbf{Śr.\ stopień} \\
\hline
Gdańsk Południe & 5\,159   & 11\,745  & 4{,}55 \\
Gdańsk          & 24\,610  & 53\,584  & 4{,}35 \\
Londyn          & 186\,437 & 300\,862 & 3{,}23 \\
\hline
\end{tabular}
\end{table}

Po wstawieniu lokalizacji usług (procedura opisana w~podrozdziale~\ref{sec:dane:wstawianie})
rozmiary grafów rosną o~liczbę wstawionych wierzchołków i~krawędzi.
Tabela~\ref{tab:grafy_po} przedstawia rozmiary grafów po wstawieniu usług
wraz z~liczbą lokalizacji w~poszczególnych kategoriach.

\begin{table}[ht]
\centering
\caption{Rozmiary grafów po wstawieniu lokalizacji usług.}
\label{tab:grafy_po}
\begin{tabular}{l r r r}
\hline
\textbf{Obszar} & $|V|$ & $|E|$ & \textbf{Usługi} \\
\hline
Gdańsk Południe & 6\,106   & 12\,692  & 947 \\
Gdańsk          & 32\,004  & 60\,978  & 7\,394 \\
Londyn          & 270\,454 & 384\,879 & 84\,017 \\
\hline
\end{tabular}
\end{table}

\begin{table}[ht]
\centering
\caption{Liczba lokalizacji usług w~poszczególnych kategoriach NEXI.}
\label{tab:uslugi_kategorie}
\begin{tabular}{l r r r}
\hline
\textbf{Kategoria} & \textbf{Gd.\ Płd.} & \textbf{Gdańsk} & \textbf{Londyn} \\
\hline
education     & 61   & 406   & 4\,378 \\
entertainment & 17   & 183   & 2\,910 \\
grocery       & 221  & 1\,377 & 10\,507 \\
health        & 77   & 557   & 4\,821 \\
posts\_banks  & 73   & 512   & 3\,249 \\
parks         & 22   & 219   & 3\,320 \\
sustenance    & 99   & 1\,207 & 26\,370 \\
shops         & 377  & 2\,933 & 28\,462 \\
\hline
\textbf{Łącznie} & \textbf{947} & \textbf{7\,394} & \textbf{84\,017} \\
\hline
\end{tabular}
\end{table}

Uzasadnienie wyboru tych trzech skali (lokalna, miejska, metropolia)
zostało przedstawione we Wprowadzeniu (podrozdział~\ref{sec:intro:metodyka:obszary}).
